\documentclass[a4paper,12pt]{article}
\usepackage[utf8]{inputenc}
\usepackage[ngerman]{babel}
\usepackage{amsmath}
\usepackage[pdftex]{graphicx}
\usepackage{geometry}
\geometry{a4paper, top=15mm, left=15mm, right=15mm, bottom=15mm,
	headsep=10mm, footskip=12mm}
\usepackage{parskip}
\usepackage{href-ul}
\hypersetup{
	colorlinks=true,
	linkcolor=blue,
	urlcolor=blue}
	



\begin{document}
	
Übungsblatt: Einfache elektrische Schaltungen 

	\section*{Grundwissen – zum Ankreuzen oder Ausfüllen}
	
	\begin{itemize}
		\item \textbf{Ein geschlossener Stromkreis} ist notwendig, damit eine Lampe \underline{\hspace{3cm}}.
		
		\item Der Strom fließt vom \underline{\hspace{2.5cm}}-Pol der Batterie zum \underline{\hspace{2.5cm}}-Pol.
		
		\item In einem Stromkreis braucht man mindestens:
		\begin{itemize}
			\item ( ) eine Steckdose
			\item ( ) eine Batterie oder Spannungsquelle
			\item ( ) einen Verbraucher (z. B. Lampe)
			\item ( ) eine elektrische Zahnbürste
			\item ( ) Kabel oder Leitungen
		\end{itemize}
		
		\item Zeichne das Symbol für:
		\begin{itemize}
			\item Batterie: \underline{\hspace{5cm}}
			\item Lampe: \underline{\hspace{5cm}}
			\item Schalter: \underline{\hspace{5cm}}
		\end{itemize}
	\end{itemize}
	
	\vspace{1em}
	\hrule
	\vspace{1em}
	
	\section*{ Stromkreis verstehen }
	
	Aurika bastelt einen einfachen Stromkreis mit einer Batterie, einem Kabel, einem Schalter und einer Glühbirne.  
	Sie schließt den Schalter – und siehe da: die Lampe leuchtet!
	
	\begin{itemize}
		\item a) Warum muss der Stromkreis „geschlossen“ sein?
		
		\vspace{1em}
		\underline{\hspace{12cm}}
		
		\item b) Jetzt hat Aurika zwei gleiche Glühbirnen und eine Stromquelle
		Welche Möglichkeiten gibt es für die  Schaltung? Zeichne zwei mögliche Schaltpläne dafür!
		Was muss man jeweils beachten ?
		
\newpage		
		\item c) Kreuze an:  
		In einer {\bf Reihenschaltung } zweier Lampen...
		
		\begin{itemize}
			\item( ) kann man beide Lampen unabhängig voneinander aus- und einschalten.
			\item( ) gehen beide Lampen aus, wenn eine durchbrennt.
			\item( ) fließt derselbe Strom durch alle Lampen, auch wenn sie verschieden sind.
			\item( ) wird es heller, wenn man mehr Lampen einbaut.
	        \item( ) müssen beide Lampen die gleiche Betriebsstromstärke in Ampere haben.
	        \item( ) müssen beide Lampen die gleiche Betriebsspannung in Volt haben.
	        \item( ) müssen beide Lampen die gleiche Leistung  in Watt haben.
	        \item( ) muss die Summe der Betriebsspannungen der Lampen gleich der Spannung der Stromquelle sein.
		\end{itemize}
		
		\item d) Kreuze an:  
		In einer {\bf Parallelschaltung } zweier Lampen...
		
		\begin{itemize}
			\item( ) kann man beide Lampen unabhängig voneinander aus- und einschalten.
			\item( ) gehen beide Lampen aus, wenn eine durchbrennt.
			\item( ) fließt derselbe Strom durch alle Lampen, auch wenn sie verschieden sind.
			\item( ) wird es heller, wenn man mehr Lampen einbaut.
			 \item( ) müssen beide Lampen die gleiche Betriebsstromstärke in Ampere haben.
			\item( ) müssen beide Lampen die gleiche Betriebsspannung in Volt haben.
			\item( ) müssen beide Lampen die gleiche Leistung  in Watt haben.
			\item( ) muss die Summe der Betriebsspannungen der Lampen gleich der Spannung der Stromquelle sein.
		\end{itemize}
		
 \item e) Jetzt willst Du die Lampen in den Stromkreisen an- und ausschalten können. 
      
      {\bf Reihenschaltung}: Erkläre, warum hier nur ein Schalter genügt und zeichne diesen in den Stromkreis ein.
      Was passiert, wenn der Schalter betätigt wird?
       
      {\bf Parallelschaltung}: Erkläre, warum hier zwei Schalter Sinn machen. Zeichne diese in den Stromkreis ein.
      Was passiert, wenn ein Schalter davon betätigt wird? Und was, wenn beide betätigt werden?
		
	\item f) Werden im Haushalt die elektrischen Geräte, wie z. B. Computer, Staubsauger usw. parallel oder in Reihe an das 230 V Netz geschaltet?
	\end{itemize}
	

\fbox{
	\begin{minipage}{0.5\textwidth}
		Zur Lösung bitte \href{https://www.okuyakl.de/phys/p7escL129/ll129.pdf}{hier klicken} oder den QR-Code scannen.\\
		Weitere Arbeitsblätter gibt es unter 
		
		\href{https://www.okuyakl.de}{www.okuyakl.de}
	\end{minipage}
	\hfill
	\begin{minipage}{0.4\textwidth}
		\includegraphics[width=1.5 cm]{../../viecher/zwe03}
		\includegraphics[width=3 cm]{qr129}
		\includegraphics[width=2 cm]{../../viecher/afanticon1}
		
\end{minipage}}
\end{document}
